\section{Introduction}\label{I:sec:intro}

A well-known theorem of Moreira, Richter and Robertson~\cite{moreira2019proof}, answering a question of
Erd\H{o}s, says that every set $A\subset\N$ of positive upper Banach density contains a sumset
$B+C=\{b+c:b\in B,\ c\in C\}$ with $B$ and $C$ infinite. In the survey~\cite{kra2023problems},
Kra, Moreira, Richter and Robertson collected a large number of open problems around this theorem:
which other infinite configurations must occur in sets of positive density, which density
hypotheses are needed, and what happens in the primes, in random sets, in multiplicative settings and
in other groups.

In this paper we answer twenty of the questions posed in~\cite{kra2023problems}. Our main result
concerns sums of more than two infinite sets in abelian groups.

\begin{theorem*}[\cref{K:thm:main}]
Let $\Gamma$ be a countably infinite abelian group and let $k\ge1$. Every set $A\subset\Gamma$ of
positive upper Banach density contains a sumset $B_1+\dots+B_k$, where $B_1,\dots,B_k\subset\Gamma$
are infinite.
\end{theorem*}

This answers Question~5.14 of~\cite{kra2023problems}. The proof does not use the Host--Kra structure
theory or the theory of nilspaces: it extends the short proof of the case $k=2$
by Kra, Moreira, Richter and Robertson~\cite{kra2026short} to an arbitrary number of summands,
building the required configuration one level at a time and keeping track only of the supports of
suitable measures on the Kronecker factor.

The remaining results are summarised in \cref{I:tab:results}. In addition, in \cref{KR:sec:main} we answer a question of Kousek and Radi\'c~\cite{kousek2024infinite}
on lower Banach density: the thresholds for $B+B+t$ and for $B+B$ are exactly $1/2$ and $3/4$. Several of them are short, and some rely
in an essential way on recent work of other authors; we indicate this in each case.

\subsection*{Conventions}
We use the numbering of the second arXiv version of~\cite{kra2023problems}. Some questions in the
survey are posed in the running text without a number. We refer to them by the subsection in which
they appear followed by a letter, in order of appearance: for instance, Question~\S3.2(a) is the
first unnumbered question in Subsection~3.2 of~\cite{kra2023problems}. Before answering each
question we restate it in our notation.

Throughout, $\N=\{1,2,3,\dots\}$, $\P$ denotes the set of primes, and for $A\subset\N$
\[
\udens(A)=\limsup_{N\to\infty}\frac{|A\cap[N]|}{N},\qquad
\ubd(A)=\lim_{N\to\infty}\sup_{M\ge0}\frac{|A\cap[M+1,M+N]|}{N},\qquad
\udlog(A)=\limsup_{N\to\infty}\frac{1}{\log N}\sum_{n\in A\cap[N]}\frac1n,
\]
where $[N]=\{1,\dots,N\}$. Lower densities $\ldens$ and natural and logarithmic densities $\dens$,
$\dlog$ are defined in the same way. In a countable abelian group, or more generally a countable
amenable group, upper Banach density $\ubd$ is the supremum of upper densities along \Fol{} sequences.
A set is \emph{thick} if it contains arbitrarily long intervals (in a group: translates of every
finite set). For sets $B,C$ we write $B+C=\{b+c:b\in B,c\in C\}$, $B+t=\{b+t:b\in B\}$ and
$mB=\{mb:b\in B\}$.

\subsection*{Questions whose answer depends on the formulation}
For several questions the answer we give concerns the question exactly as it is stated in the survey,
and the literal statement differs from what the authors of the survey probably intended. We include
these answers because the difference is not always obvious and the literal statements have been
repeated in the literature, but in each case we point out the issue explicitly, give the natural
corrected formulation, and say whether it remains open. This applies to
Question~3.20 (where $m=1$ is coprime to every $q$), Question~3.23 (where $B+mB$ contains the
diagonal terms $(m+1)b$), Question~\S3.2(a) (where the density condition can be read in two ways; we
answer both readings), Question~5.6, Question~5.23 and Question~\S5.4(b) (where the fibre conditions
constrain only one coordinate). The negative answer to Question~\S3.3(e) follows directly from known
counterexamples (one from~\cite{kra2023problems} itself and one due to
Ackelsberg~\cite{ackelsberg2024counterexamples}); we record it for completeness. Similarly, our
answer to Question~3.24 is preceded by an example of Kousek and Radi\'c~\cite{kousek2024infinite},
which, as we check, already answers that question.

\begin{table}[ht]
\centering
\small
\renewcommand{\arraystretch}{1.25}
\begin{tabular}{@{}p{0.17\textwidth}p{0.58\textwidth}p{0.14\textwidth}@{}}
\hline
Question & Answer & Section\\
\hline
5.14 & Yes: $B_1+\dots+B_k$ in every countably infinite abelian group & \ref{K:sec:main}\\
5.16, \S5.3(a) & Exactly which sets of translates suffice when $[G:2G]<\infty$ & \ref{G:sec:translates}\\
5.10 & No: fibrewise lower density does not suffice in $\N^2$ & \ref{G:sec:q510}\\
5.23, \S5.4(b) & No, as stated; a corrected version is open & \ref{G:sec:q523}\\
3.20 & No (for $B+mB$ as in the survey), for two and for three colours & \ref{IA:sec:q320}\\
3.23 & No, as stated & \ref{IA:sec:q323}\\
3.24 & No (already implied by Kousek--Radi\'c) & \ref{IA:sec:q324}\\
\S3.2(a) & No, for both readings & \ref{IA:sec:u32a}\\
\S3.3(a) & Ordered polynomial sums: exactly the linear case & \ref{IA:sec:u33a}\\
\S3.3(e) & No (direct consequence of known examples) & \ref{IA:sec:u33e}\\
3.31, \S3.6(a) & Yes, in the strongest $k$-fold form & \ref{IB:sec:q331}\\
4.13 & Yes iff $\gamma\notin\Z$ (for almost every $x$) & \ref{IB:sec:q413}\\
\S4.4(a) & Random sets: criterion for monotone $p_n$; no Borel criterion in general & \ref{IB:sec:u44a}\\
4.5 & Yes, unconditionally & \ref{IB:sec:c45}\\
5.4 & No & \ref{SP:sec:q54}\\
5.6 & No, as stated & \ref{SP:sec:q56}\\
\S5.1(a) & Yes for both examples & \ref{SP:sec:u51a}\\
\hline
\end{tabular}
\caption{Questions from~\cite{kra2023problems} answered in this paper.}
\label{I:tab:results}
\end{table}

\subsection*{Formalization}
Many of the results have been formalized in Lean~4 with Mathlib, among them the answers to
Questions~3.20, 3.23, 3.24, \S3.2(a), \S3.3(e), \S4.4(a), 5.4, 5.6, 5.10, 5.23 and \S5.1(a), the
integer case of Question~4.13, and the result of \cref{KR:sec:main}. A few further results are formalized conditionally on the
published theorems they rely on, which are stated as explicit hypotheses. The formalization is
available at \url{https://github.com/daizisheng/kmrr25}.

\subsection*{Organisation}
\Cref{K:sec:main} contains the proof of the main theorem. The rest of the paper is organised by the
setting: further results on groups and on $\N^2$, results on sumsets in the integers, and results on
configurations mixing sums and products. The sections can be read independently.
